% !TeX root = ../main.tex
\section{Related Work}
\label{sec:related-work}

Hardware-aware training (HAT) mitigates analog non-idealities by injecting hardware perturbations into the forward path during optimization \citep{joshi2020accurate}. It is related to quantization-aware training \citep{choi2019pact} and to domain randomization in sim-to-real transfer \citep{tobin2017domain}, but CIM introduces a distinct transfer problem: device-to-device (D2D) mismatch is spatially structured and fixed for a fabricated instance, whereas cycle-to-cycle (C2C) noise is stochastic per read. Training against one fixed mismatch field can therefore produce an instance-specific solution that does not transfer. Ensemble HAT applies the domain-randomization idea at the hardware-instance level by resampling the static D2D mask across epochs.

The need for this distinction is especially visible in organic neuromorphic and optoelectronic devices. These devices offer optical sensitivity, multilevel conductance tuning, and flexible-substrate compatibility for neuromorphic hardware \citep{xu2025emerging,photonics2025organicreview,beller2024organicneurons,ji2025singleoectneuron}. Recent demonstrations report multilevel memory \citep{guo2024organic,jung2024organicfilaments}, long retention tails \citep{zeng2023organicmemristor}, and integrated optoelectronic arrays \citep{gebregiorgis2023organiccim,liu2026opect,zhang2025mooptoelectronic,cui2025multimode}. The system question is now becoming more explicit: active-matrix addressing, spatial optical response, crosstalk, and temperature stability are increasingly treated as deployment constraints rather than only as materials observations \citep{visionarch2023crosstalk,amspa2024insensor,jang2023insensor,fuller2020tempresilient,guo2024hightemp}. This motivates network-level tools that can substitute device profiles and rank failure modes before a full array is available.

Existing CIM simulators provide the foundation for such studies. DNN+NeuroSim established an effective template for connecting device assumptions, peripheral-circuit costs, and neural-network accuracy \citep{peng2020dnnneurosim}. MemTorch embeds memristive non-idealities into PyTorch workflows for end-to-end evaluation \citep{lammie2022memtorch}, AIHWKit provides analog HAT and inference simulation in a practical mixed-signal software stack \citep{rasch2021aihwkit}, and CrossSim supplies a GPU-accelerated crossbar-accuracy workflow with parasitic resistance, ADC, drift, and lookup-table-based training models \citep{crosssim2024}. These mature simulators remain the correct tools for circuit-aware studies, but they do not expose every profile-level effect of interest, such as inverse-gamma optoelectronic photoresponse or organic-specific double-exponential retention, as first-class device primitives. Our behavioral profile-driven layer is therefore positioned as a deployment-risk ranking extension, not as a replacement for established CIM toolkits.
